% Part: normal-modal-logic
% Chapter: filtrations
% Section: S5-fmp

\documentclass[../../../include/open-logic-section]{subfiles}

\begin{document}

\olfileid{nml}{fil}{fmp}

\olsection{\Log{K}, \Log{S5}లకు పరిమిత నమూనా ధర్మం ఉంది}

\begin{defn}
  మోడల్ తర్క వ్యవస్థ $\Sigma$కు \emph{పరిమిత నమూనా ధర్మం}
  ఉందంటే, ఏ \tetoken{సూత్రం}{!!a{formula}}~$!A$
  అయినా $\Sigma$కు చెందిన ఒక నమూనాలో ఏదో లోకంలో సత్యమైతే,
  $!A$ ఒక \emph{పరిమిత} $\Sigma$ నమూనాలో
  ఏదో లోకంలో సత్యమవుతుంది.
\end{defn}

\begin{prop}\ollabel{prop:K-fmp}
  \Log{K}కు పరిమిత నమూనా ధర్మం ఉంది.
\end{prop}

\begin{proof}
  \Log{K} చెల్లుబాటయ్యే
  \tetoken{సూత్రాల}{!!{formula}s} సమితి. అంటే ప్రతి
  నమూనా~\Log{K} నమూనాయే. ఇప్పుడు
  \olref[fil]{thm:filtrations} ప్రకారం,
  $\mSat{M}{!A}[w]$ అయితే,~$!A$కు చెందిన
  \tetoken{ఉపసూత్రాల}{!!{formula}s} సమితి $\Gamma$
  ద్వారా~$\mModel{M}$కు చేసే ఏ వడపోతలోనైనా
  $\mSat{M^*}{!A}[{[w]}]$.
  ఏ \tetoken{సూత్రానికైనా}{!!{formula}} పరిమిత
  సంఖ్యలోనే \tetoken{ఉపసూత్రాలు}{!!{formula}s}
  ఉంటాయి; కనుక $\Gamma$ పరిమితం.
  \olref[fin]{prop:filt-are-finite} ప్రకారం
  $\card{W^*} \le 2^n$; ఇక్కడ $n$ అనేది
  $\Gamma$లోని \tetoken{సూత్రాల}{!!{formula}s} సంఖ్య.
  \Log{K} నమూనాలపై ఏ ఆంక్షా విధించదు కనుక
  $\mModel{M^*}$ ఒక \Log{K}-నమూనా.
\end{proof}

వడపోతల ద్వారా తర్కం~\Log{L}కు పరిమిత నమూనా
ధర్మం ఉందని చూపాలంటే, \Log{L}-నమూనాకు మనం ఎంచుకునే
వడపోత కూడా \Log{L}-నమూనా కావాలి. ప్రతి వడపోతా
\Log{L}-నమూనాను ఇవ్వదు కనుక దీనికి తరచూ చాలా పని అవసరం.
సార్వత్రిక నమూనాల వడపోతలు మాత్రం సార్వత్రికంగానే ఉంటాయి.

\begin{prop}\ollabel{prop:univ-fin}
  $\mClass{U}$ సార్వత్రిక నమూనాల వర్గం
  (\olref[frd][es5]{prop:S5=univ} చూడండి),
  $\mClass{U}_\mathrm{Fin}$ అన్ని పరిమిత సార్వత్రిక
  నమూనాల వర్గం అనుకుందాం. అప్పుడు ఏ
  \tetoken{సూత్రం}{!!{formula}}~$!A$ అయినా
  $\mClass{U}$లో చెల్లుబాటయ్యేది
  $\mClass{U}_\mathrm{Fin}$లో చెల్లుబాటైనప్పుడు, అప్పుడే మాత్రమే.
\end{prop}

\begin{proof}
  పరిమిత సార్వత్రిక నమూనాలు కూడా సార్వత్రిక నమూనాలే.
  కనుక ఎడమ నుంచి కుడికి దిశ నేరుగా వస్తుంది.
  తిరుగు దిశలో, సార్వత్రిక నమూనా
  $\mModel{M}$లోని ఏదో లోకం $w$లో~$!A$
  అసత్యమని అనుకుందాం. $\Gamma$లో $!A$తోపాటు
  దాని అన్ని ఉపసూత్రాలను చేర్చుదాం; స్పష్టంగా
  $\Gamma$ పరిమితం. $\mModel{M}$కు ఒక వడపోత
  $\mModel{M^*}$ను తీసుకుందాం. అప్పుడు
  \olref[fin]{prop:filt-are-finite} ప్రకారం
  $\mModel{M^*}$ పరిమితం; అలాగే
  \olref[fil]{thm:filtrations} ప్రకారం
  $\mModel{M^*}$లో $[w]$ అనే లోకంలో $!A$ అసత్యం.
  $\mModel{M^*}$ కూడా సార్వత్రికమని గమనిస్తే
  సరిపోతుంది: ఏ $u$, $v$లనైనా తీసుకుంటే,
  పరికల్పన ప్రకారం $Ruv$; అందువల్ల
  \olref[fil]{defn:filtration}\olref[fil]{defn:filtration-R}
  ప్రకారం $R^*[u][v]$ కూడా.
\end{proof}

\begin{cor}\ollabel{cor:S5fmp}
  \Log{S5}కు పరిమిత నమూనా ధర్మం ఉంది.
\end{cor}

\begin{proof}
  \olref[frd][es5]{prop:S5=univ} ప్రకారం,
  $!A$ ఒక స్వావర్తన, యూక్లిడియన్ నమూనాలోని
  ఏదో లోకంలో సత్యమైతే, అది ఒక సార్వత్రిక
  నమూనాలోని ఏదో లోకంలో సత్యం. \olref{prop:univ-fin}లోని
ఫలితం ప్రకారం, అది ఒక పరిమిత
సార్వత్రిక నమూనాలోని ఏదో లోకంలో సత్యం.
(ఇక్కడ ఆ పరిమిత నమూనా,~$!A$కు చెందిన
  \tetoken{ఉపసూత్రాల}{!!{formula}s} సమితి ద్వారా
  ఆ నమూనాకు చేసే వడపోతలో). ప్రతి సార్వత్రిక
  నమూనాకు స్వావర్తన, యూక్లిడియన్ ధర్మాలు రెండూ ఉంటాయి కనుక,
  $!A$ ఒక పరిమిత స్వావర్తన యూక్లిడియన్ నమూనాలోని
  ఏదో లోకంలో సత్యం.
\end{proof}

\begin{prob}
  సీరియల్ లేదా స్వావర్తన నమూనాకు చేసే ఏ వడపోతైనా
  వరుసగా సీరియల్ లేదా స్వావర్తన నమూనానే అని చూపండి.
\end{prob}

\begin{prob}
  సౌష్ఠవ (సంక్రామక, యూక్లిడియన్) నమూనాకు
  సౌష్ఠవం లేని (సంక్రామకత్వం లేని, యూక్లిడియన్
  కాని) వడపోతను కనుగొనండి.
\end{prob}

\end{document}
